%% kapa: exam=1 round=36 year=2025 period=1 form=A subject=economics subject_ko=경제학원론 seq=4 number=44
%% source: Q-Net 감정평가사 기출문제 (cst00309) artlSeq=5231525 file=raw/1008/5231525/2025년 제36회 감정평가사 1차시험 1교시 문제지 원본.pdf
%% conversion: kordoc 4.24.51 -> tools/gen1.py (auto); review=proofread-sonnet; flags=none
\begin{question}{44}
생산가능곡선에 관한 설명으로 옳은 것을 모두 고른 것은?
\begin{framedtext}
ㄱ. 곡선의 외부에 있는 점은 비효율적인 생산점이고, 내부에 있는 점은 실현 불가능한 생산점이다. \par
ㄴ. 곡선이 원점에 대해 오목하면 한 재화의 생산을 늘릴수록 기회비용이 증가한다. \par
ㄷ. 곡선이 직선이면 한 재화의 생산을 늘릴수록 기회비용이 감소한다.
\end{framedtext}
\begin{choices}
  \item ㄱ
  \item ㄴ
  \item ㄱ, ㄷ
  \item ㄴ, ㄷ
  \item ㄱ, ㄴ, ㄷ
\end{choices}
\end{question}
\answer{2}
